\begin{problem}[40.6]
  Find the degree and a basis for the field extension $\Q(\sqrt{2} + \sqrt{3})$ over $\Q$.
\end{problem}

\begin{solution}
  The basis for the field extension $\Q(\sqrt{2} + \sqrt{3})$ over $\Q$ is $\{1, \sqrt{2} + \sqrt{3}, (\sqrt{2} + \sqrt{3})^2, (\sqrt{2} + \sqrt{3})^3\}$. The degree of the field extension is 4, since the minimal polynomial of $\sqrt{2} + \sqrt{3}$ over $\Q$ is given by $f(x) = x^4 - 10x + 1$. This polynomial is irreducible over $\Q$ by Eisenstein's criterion with $p = 5$, and thus the degree of the extension is equal to the degree of the minimal polynomial, which is 4.
\end{solution}

\begin{problem}[40.8]
  Find the degree and a basis for the field extension $\Q(\sqrt{2}, \sqrt[3]{5})$ over $\Q$.
\end{problem}

\begin{solution}
  The basis for the field extension $\Q(\sqrt{2}, \sqrt[3]{5})$ over $\Q$ is $\{1, \sqrt{2}, \sqrt[3]{5}, \sqrt{2}\sqrt[3]{5}, (\sqrt[3]{5})^2, \sqrt{2}(\sqrt[3]{5})^2\}$. The degree of the field extension is 6, since the minimal polynomial of $\sqrt{2}$ over $\Q$ is $x^2 - 2$, which has degree 2, and the minimal polynomial of $\sqrt[3]{5}$ over $\Q$ is $x^3 - 5$, which has degree 3. Since these two extensions are linearly disjoint, the degree of the combined extension is the product of the degrees, which is $2 \cdot 3 = 6$.
\end{solution}

\begin{problem}[40.10]
  Find the degree and a basis for the field extension $\Q(\sqrt{2}, \sqrt{6})$ over $\Q(\sqrt{3})$.
\end{problem}

\begin{solution}
  Note that $\sqrt{6}/\sqrt{2} = \sqrt{3}$. Thus, $\Q(\sqrt{2}, \sqrt{6}) = \Q(\sqrt{2}, \sqrt{3})$. This problem is now reduced to finding the basis and degree of the field extension $\Q(\sqrt{2}, \sqrt{3})$ over $\Q(\sqrt{3})$. The basis for this extension is $\{1, \sqrt{2}\}$, since $\sqrt{2}$ is not in $\Q(\sqrt{3})$. The degree of the field extension is 2, since the minimal polynomial of $\sqrt{2}$ over $\Q(\sqrt{3})$ is $x^2 - 2$, which has degree 2. Thus, the degree of the field extension $\Q(\sqrt{2}, \sqrt{6})$ over $\Q(\sqrt{3})$ is 2, and a basis for this extension is $\{1, \sqrt{2}\}$.
\end{solution}

\begin{problem}[40.18]
  Show by an example that for a proper extension field $E$ of a field $F$, the algebraic closure of $F$ in $E$ need not be algebraically closed.
\end{problem}

\begin{solution}
  Consider the field extension $\Q(\sqrt{2})$ over $\Q$. The algebraic closure of $\Q$ in $\Q(\sqrt{2})$ is $\Q$, which is not algebraically closed since it does not contain all the roots of polynomials with coefficients in $\Q$. For example, the polynomial $x^2 - 2$ has a root $\sqrt{2}$ in $\Q(\sqrt{2})$, but it does not have any roots in $\Q$. Thus, the algebraic closure of $\Q$ in $\Q(\sqrt{2})$ is not algebraically closed.
\end{solution}

\begin{problem}[40.23]
  Show that if $E$ is a finite extension of a field $F$ and $[E : F]$ is a prime number, then $E$ is a simple extension of $F$ and, indeed, $E = F(\alpha)$ for every $\alpha \in E$ not in $F$.
\end{problem}

\begin{solution}
  Since $\alpha$ is algebraic over $F$, there exists a minimal polynomial $f(x) \in F[x]$ such that $f(\alpha) = 0$. The degree of this minimal polynomial divides the degree of the extension, which is $p$. Since $p$ is prime, the degree of the minimal polynomial must be either 1 or $p$. If the degree is 1, then $\alpha \in F$, which contradicts our assumption. Therefore, the degree of the minimal polynomial must be $p$, and thus $[F(\alpha) : F] = p$.

  Since $[E : F] = p$ and $[F(\alpha) : F] = p$, we have that $E = F(\alpha)$, as there can be no intermediate fields between $F$ and $E$. Therefore, every element $\alpha \in E$ that is not in $F$ generates the entire extension, and we conclude that $E$ is a simple extension of $F$.
\end{solution}

\begin{problem}[40.24]
  Prove that $x^2 - 3$ is irreducible over $\Q(\sqrt[3]{2})$.
\end{problem}

\begin{solution}
  We know that $[\Q(\sqrt{3}) : \Q] = 2$ and $[\Q(\sqrt[3]{2} : \Q] = 3$. Now, suppose for contradiction, that $\sqrt{3} \in \Q(\sqrt[3]{2})$. Then, $\Q(\sqrt{3}) \subseteq \Q(\sqrt[3]{2})$. Then, we have
  \[%
    [\Q(\sqrt[3]{2}) : \Q] = [\Q(\sqrt[3]{2}) : \Q(\sqrt{3})][\Q(\sqrt{3}) : \Q]
  .\]%
  So, we have
  \[%
    3 = [\Q(\sqrt[3]{2}) : \Q(\sqrt{3})] \cdot 2
  ,\]%
  which is a contradiction since the left-hand side is odd and the right-hand side is even. Thus, $\sqrt{3} \notin \Q(\sqrt[3]{2})$, and therefore $x^2 - 3$ is irreducible over $\Q(\sqrt[3]{2})$.
\end{solution}

\begin{problem}[40.25]
  What degree field extensions can we obtain by successively adjoining to a field $F$ a square root of an element of $F$ not a square in $F$, then a square root of some nonsquare in this new field, and so on? Argue from this that a zero of $x^{14} - 3x^2 + 12$ over $\Q$ can never be expressed as a rational function of square roots of rational functions of square roots, and so on, of elements of $\Q$.
\end{problem}

\begin{solution}
  By successively adjoining square roots of nonsquares, we can only obtain field extensions of degree $2^n$ for some nonnegative integer $n$. This is because each time we adjoin a square root, we either get a degree 1 extension (if the element is already a square) or a degree 2 extension (if the element is not a square). Therefore, the degree of the resulting field extension after $n$ steps will be at most $2^n$.

  Now, consider the polynomial $x^{14} - 3x^2 + 12$. The degree of this polynomial is 14, which is not a power of 2. Therefore, any field extension that contains a zero of this polynomial must have degree that divides 14. Since the only powers of 2 that divide 14 are 1 and 2, it follows that any field extension containing a zero of this polynomial cannot be obtained by successively adjoining square roots as described. Thus, a zero of $x^{14} - 3x^2 + 12$ over $\Q$ can never be expressed as a rational function of square roots of rational functions of square roots, and so on, of elements of $\Q$.
\end{solution}

\begin{problem}[40.28]
  Generalizing Exercise 27, show that if $\sqrt{a} + \sqrt{b} \ne 0$, then $\Q(\sqrt{a} + \sqrt{b}) = \Q(\sqrt{a}, \sqrt{b})$ for all $a$ and $b$ in $\Q$. [\textit{Hint:} Compute $(a - b)/(\sqrt{a} + \sqrt{b})$.]
\end{problem}

\begin{solution}
  Computing $(a - b)/(\sqrt{a} + \sqrt{b})$, we have
  \[%
    \frac{a - b}{\sqrt{a} + \sqrt{b}} = \frac{(a - b)(\sqrt{a} - \sqrt{b})}{(\sqrt{a} + \sqrt{b})(\sqrt{a} - \sqrt{b})} = \frac{(a - b)(\sqrt{a} - \sqrt{b})}{a - b} = \sqrt{a} - \sqrt{b}
  .\]%
  Thus, we have $\sqrt{a} - \sqrt{b} \in \Q(\sqrt{a} + \sqrt{b})$. Since $\sqrt{a} + \sqrt{b} \in \Q(\sqrt{a} + \sqrt{b})$, we have $\sqrt{a} = (\sqrt{a} + \sqrt{b} + \sqrt{a} - \sqrt{b})/2$ and $\sqrt{b} = (\sqrt{a} + \sqrt{b} - (\sqrt{a} - \sqrt{b}))/2$ are also in $\Q(\sqrt{a} + \sqrt{b})$. Thus, we have $\Q(\sqrt{a}, \sqrt{b}) \subseteq \Q(\sqrt{a} + \sqrt{b})$. The other inclusion is trivial since $\Q(\sqrt{a} + \sqrt{b})$ is generated by $\sqrt{a}$ and $\sqrt{b}$, which are in $\Q(\sqrt{a}, \sqrt{b})$. Therefore, we conclude that $\Q(\sqrt{a} + \sqrt{b}) = \Q(\sqrt{a}, \sqrt{b})$.
\end{solution}

\begin{problem}[40.29]
  Let $E$ be a finite extension of a field $F$, and let $p(x) \in F[x]$ be irreducible over $F$ and have degree that is not a divisor of $[E : F]$. Show that $p(x)$ has no zeros in $E$.
\end{problem}

\begin{solution}
  Suppose for contradiction that $p(x)$ has a zero $\alpha$ in $E$. Then, we have the field extension $F \subseteq F(\alpha) \subseteq E$. Since $p(x)$ is irreducible over $F$, the degree of the extension $[F(\alpha) : F]$ is equal to the degree of $p(x)$, which we denote as $n$. Thus, we have
  \[%
    [E : F] = [E : F(\alpha)][F(\alpha) : F] = [E : F(\alpha)]n
  .\]%
  This implies that $n$ divides $[E : F]$, which contradicts our assumption. Therefore, we conclude that $p(x)$ has no zeros in $E$.
\end{solution}

\begin{problem}[40.30]
  Let $E$ be an extension field of $F$. Let $\alpha \in E$ be algebraic of odd degree over $F$. Show that $\alpha^2$ is algebraic of odd degree over $F$, and $F(\alpha) = F(\alpha^2)$.
\end{problem}

\begin{solution}
  We first show that if $\alpha$ is algebraic of odd degree over $F$, then $\alpha^2$ is also algebraic of odd degree over $F$. Since $\alpha^2$ is a power of $\alpha$, we have the inclusion property $F \subseteq F(\alpha^2) \subseteq F(\alpha)$. Thus, we have
  \[%
    [F(\alpha) : F] = [F(\alpha) : F(\alpha^2)][F(\alpha^2) : F]
  .\]%
  We are given that $[F(\alpha) : F] = n$, for some odd integer $n$. Consider the polynomial $f(x) = x^2 - \alpha^2 \in F(\alpha^2)[x]$. Since $\alpha$ is a root of $f(x)$, the minimal polynomial of $\alpha$ must divide $x^2 - \alpha^2$. Thus, we have $[F(\alpha) : F(\alpha^2)]$ can only be 1 or 2. If it is 1, then $F(\alpha) = F(\alpha^2)$, and we are done. If it is 2, then we have
  \[%
    n = [F(\alpha) : F] = [F(\alpha) : F(\alpha^2)][F(\alpha^2) : F] = 2[F(\alpha^2) : F]
  .\]%
  This is, however, a contradiction since $n$ is odd. Thus, we must have $[F(\alpha) : F(\alpha^2)] = 1$, and therefore $F(\alpha) = F(\alpha^2)$.
\end{solution}
